% =========================================================================
% PAPER 2 -- submission manuscript (LaTeX)
% Degree-Independent Control Impact in the Human Structural Connectome:
% Graph-Geometric Associations and Degree-Preserving Null Arbitration
% =========================================================================
% Scientific source of truth: frozen repository artifacts (v2.0.1).
% Every number is quoted from audited frozen artifacts; the numeric layer
% is independently verified by 13_MANUSCRIPT/verify_stage24.py (V01-V12,
% 12/12 PASS, deep V09 over 1,200 raw null records).
\documentclass[11pt,a4paper]{article}
\usepackage[margin=2.5cm]{geometry}
\usepackage[T1]{fontenc}
\usepackage[utf8]{inputenc}
\usepackage{amsmath,amssymb}
\usepackage{graphicx}
\usepackage{booktabs}
\usepackage{caption}
\usepackage{hyperref}
\usepackage{orcidlink} % may be removed if unavailable
\hypersetup{colorlinks=true,linkcolor=blue,citecolor=blue,urlcolor=blue}

\newcommand{\del}{\delta}

\title{\textbf{Degree-Independent Control Impact in the Human Structural
Connectome: Graph-Geometric Associations and Degree-Preserving Null
Arbitration}}

\author{Harsha Vardhan Malipeddi\thanks{Independent Researcher.
Corresponding author. ORCID/email: [author to supply before submission].}}
\date{Frozen release v2.0.1 --- September 30, 2026}

\begin{document}
\maketitle

\noindent\textbf{Funding:} The author received no specific funding for this work.\\
\textbf{Competing interests:} The author declares no competing interests.\\
\textbf{Ethics:} This study used previously acquired, de-identified data
(AOMIC ID1000); no new human data were collected.\\
\textbf{Data \& code availability:} see Section~\ref{sec:avail}.\\
\textbf{Author contributions:} H.V.M.\ conceived the study, developed the
analysis framework, performed all computational analyses, validation and
verification, and wrote the manuscript.

\begin{abstract}
\noindent\textbf{Question.} After accounting for connectivity
(degree/strength), which graph-geometric properties explain node-level
variation in removal-based control impact (CIS) in the human structural
connectome --- and do any nodes qualify as operationally defined ``hidden
chokepoints''?

\smallskip\noindent\textbf{Methods.} 801 QC-pass subjects (AOMIC ID1000;
4S456 parcellation; 15\% cost threshold) $\times$ 456 nodes. CIS consumed
read-only from the frozen Paper~1 release. The degree-independent residual
$R_i=\mathrm{CIS}_i-\widehat{\mathrm{CIS}}(\mathrm{degree}_i)$ is estimated
per subject with a frozen cubic-spline family (orthogonality gate: median
$|\rho(R,\mathrm{degree})|=0.0327$). Eleven network features are associated
with $R$ within-subject, arbitrated by exact degree-preserving nulls
(12 subjects $\times$ 100 nulls, seed 20270927, per-null full CIS + feature
recomputation; estimator-equivalence validated), benchmarked out-of-sample
(subject-grouped CV), and mapped onto evidence-leveled biological
annotations.

\smallskip\noindent\textbf{Results.} (1) The residual is positive in all
subjects: $\del=0.104$, positive in 801/801 subjects (95\% CI
$[0.1007,0.1073]$). (2) Network position predicts the residual
out-of-sample: ridge cv-$R^2=0.574$ (permuted-$R$ control $0.0024$;
degree-only $\approx 0$), strongest associate bridge score (median
$\rho=+0.4985$). (3) \textbf{Degree-preserving-null arbitration is negative
(Outcome C):} every feature's observed association falls inside the null
p95 band --- independently deep-verified from the 1{,}200 raw null records.
(4) Joint-criterion ``chokepoint'' candidates are rare and
threshold-sensitive (1/2/3/7 nodes at $k=1/2/5/10\%$); per-node null
arbitration is not established. (5) Top-10\%-CIS nodes under-represent
cortex (0.444 vs 0.877; $q=0.00017$) --- permutation-supported but not
spatially controlled. (6) The fly connectome's residual also fails its
degree-preserving null ($z=1.21$, $p=0.109$, frozen): two same-direction
negatives at different scales license no universal-mechanism claim.

\smallskip\noindent\textbf{Conclusion.} Degree-independent control impact
is a real, replicated, predictable-from-network-features, anatomically
non-uniform property of human connectome architecture whose fine structure
is \textbf{null-qualified}. The ``hidden chokepoint'' remains an
operational, bounded construct rather than a validated node class.

\smallskip\noindent
\textbf{Keywords:} structural connectome; network control theory;
controllability; degree confounding; null models; graph metrics;
cross-species comparison.
\end{abstract}

\section{Introduction}
Network control theory frames the brain as a dynamical system whose
structural connectome constrains its reachable dynamics
\cite{gu2015,tang2018,liu2011}. Removal-based control impact --- here
\textbf{CIS} (Control Impact Score), the fractional loss of average network
controllability upon node removal,
$\mathrm{CIS}_i=(E(G)-E(G-i))/E(G)$ with $E$ the global efficiency
\cite{latora2001} --- indexes each region's systemic importance, extending
the error/attack-tolerance view of network robustness \cite{albert2000}.
Structural connectome construction, streamline weighting and cost
thresholding follow established recipes \cite{hagmann2008,smith2015,achard2007}.

Paper~1 of this series \cite{paper1} established two facts that motivate
the present work: (i) raw CIS in the human structural connectome is
strongly degree-confounded (median $\rho=0.943$); and (ii) a small but
strictly degree-matched residual exists and replicates ($\del=0.100$;
re-derived here as $0.104$). Paper~1 also showed that anatomy does not
transfer across scales (``architecture replicates, anatomy doesn't''):
the fly visual system contributes $\approx 80\%$ of total CIS while the
human visual system contributes $\approx 2\%$.

Paper~2 asks the mechanistic follow-up, frozen before unblinding
(\texttt{01\_CONFIG/CONFIG\_FREEZE.md}; protocol + Amendments 1--3):
\begin{quote}
\textbf{H1 (primary):} degree-controlled CIS residuals are associated with
specific graph-geometric properties beyond degree/strength.
\end{quote}
Secondary hypotheses H2--H5 concern path redundancy, inter-community
participation, core/bridge positions, and multivariate explainability;
H6--H7 (cross-scale) are exploratory and are never promoted to
confirmatory claims. All confirmatory tests are pre-registered against an
Outcome A/B/C arbitration rule: \textbf{A} = association survives an exact
degree-preserving null (mechanism licensed); \textbf{B} = mixed;
\textbf{C} = no association exceeds the null band (mechanism not licensed).

\section{Data and node semantics}
\textbf{Human.} AOMIC ID1000 derived structural connectomes
\cite{snoek2021} (Zenodo 19796783, CC-BY-4.0): 801 QC-pass subjects, 4S456
parcellation, SIFT streamline counts \cite{smith2015}, 15\% cost threshold.
CIS values are consumed \textbf{read-only} from the frozen Paper~1 release.
\textbf{Fly (comparison only).} FAFB v783 frozen release artifacts
\cite{dorkenwald2024,dorkenwald2023}: individual neurons; 3{,}518 targets.
\textbf{Node semantics (mandatory).} A human network node is an atlas
\textbf{parcel} --- a macroscopic brain region, never an individual neuron.

\section{Methods}
\subsection{CIS and the degree-independent residual}
$\mathrm{CIS}_i=(E(G)-E(G-i))/E(G)$, exact recomputation. The residual is
$R_i=\mathrm{CIS}_i-\mathbb{E}[\mathrm{CIS}\,|\,\mathrm{degree}_i]$ with
the conditional expectation estimated \textbf{per subject} by OLS of raw
CIS on a df\,=\,4 natural cubic spline of degree \cite{hastie2009}. Gates
(all passed): orthogonality median $|\rho(R,\mathrm{degree})|=0.0327<0.05$;
same-pairs agreement with Paper~1's matched-pair Cliff's $\del$
\cite{cliff1993}, $\rho=0.785$. No transform, offset, or retransformation
exists in the final estimator (history: Amendments 1--2).

\subsection{Feature set and inference}
Strength, betweenness \cite{freeman1977}, closeness, path redundancy,
bridge score, participation and within-module $z$
\cite{guimera2005}, k-core \cite{seidman1983}, eigenvector centrality
\cite{bonacich1972}, PageRank \cite{page1999}, clustering
\cite{watts1998} --- 11 features + degree. Within-subject Spearman
associations (Stage 8); confirmatory within-subject standardized
regression with CR1 cluster-robust SEs \cite{mackinnon1985} clustered on
subject ($N=801$, 365{,}256 rows); 5-fold subject-grouped CV $\times$ 3
seeds for the predictive benchmark (predictive performance is \textbf{not}
a mechanism claim).

\subsection{Degree-preserving null arbitration (confirmatory)}
Maslov--Sneppen double-edge rewiring preserving the \textbf{exact degree
sequence} \cite{maslov2002}, 12 subjects $\times$ 100 nulls, seed 20270927;
per null: full CIS recomputation + full feature recomputation +
residualization via the documented 20-bin quantile-median proxy of the
frozen spline. The proxy is validated on observed data: the arbitrated
association statistic differs between proxy and spline by a median of 0.050
--- far below the null bands (0.10--0.73) --- and the binned variant yields
\emph{stronger} observed associations, so it cannot manufacture Outcome~C.

\paragraph{Null ensemble calibration and Monte Carlo precision.}
The 12 subjects are a frozen-seed random draw from the 801; their pooled
statistics are representative (bridge median $+0.466$ vs $+0.4985$
population). 100 nulls/subject $\times$ 12 follows the Paper-1 precedent
and the compute cost of exact CIS per null. With 1{,}200 records the p95
rests on the upper $\sim$60 order statistics; its Monte-Carlo error is
small relative to the band widths (per-feature null IQR of $|\rho|$
$\approx 0.03$--$0.07$), and every observed statistic sits strictly inside
its band, so the verdict is not boundary-sensitive at this ensemble size.

\subsection{Robustness, negative controls, annotation, chokepoint
definition, cross-species}
Estimator ladder + winsorization (Stage 20); negative controls NC1/NC2
(Stage 21); evidence-leveled annotation + permutation enrichment with BH
correction \cite{benjamini1995} (Stage 23); operational five-criterion
chokepoint definition (Stage 22); conservative 12-property human-vs-fly
comparison (Stage 22b). Full details: Supplementary Tables S0--S10.

\section{Results}
\subsection{The residual replicates}
$\del=0.104$ (median, 801 subjects), bootstrap 95\% CI $[0.1007,0.1073]$;
positive in \textbf{801/801} subjects; LOSO max deviation $5.4\times10^{-8}$.

\begin{table}[ht]
\centering\small
\caption{Predictive benchmark (5-fold subject-grouped CV $\times$ 3 seeds;
numpy-native). Predictive performance is not a mechanism claim.}
\begin{tabular}{lrrr}
\toprule
Model & cv-$R^2$ (mean $\pm$ sd) & cv-$R^2$ permuted & $\Delta$ vs degree-only \\
\midrule
Degree-only & $-0.00006\pm0.000002$ & $0.0004$ & --- \\
\textbf{Ridge (11 feat.)} & $\mathbf{0.5740\pm0.0001}$ & $0.0024$ & $\mathbf{+0.574}$ \\
Elastic net & $0.5614\pm0.0001$ & $0.0023$ & $+0.561$ \\
MLP & $-4.09\pm4.57$ & $-4.65$ & $-4.09$ (unstable) \\
\bottomrule
\end{tabular}
\end{table}

\subsection{Associations and controls}
All 12 features associate with $R$ at $q<10^{-27}$ (bridge $+0.4985$,
closeness $+0.2489$, redundancy $-0.2179$, participation $-0.1676$); the
8-feature CR1 model retains all at $q<10^{-7}$. NC1/NC2 chance bands are
$|\rho|\le0.0044$; all 8 features exceed both bands (margins large for
strong associates, small for the weak ones: within-module $z$ 0.014,
k-core 0.052).

\begin{table}[ht]
\centering\small
\caption{Degree-preserving null arbitration (12 subjects $\times$ 100
nulls; exact degree preservation verified per null; deep-verified by gate
V09 from the 1{,}200 raw records).}
\begin{tabular}{lrrrr}
\toprule
Feature & Obs. $|\mathrm{med}\,\rho|$ & Null med $\rho$ & Null IQR & Null $|\rho|$ p95 \\
\midrule
Bridge & 0.4985 & $-0.5957$ & 0.0502 & 0.6538 \\
Clustering & 0.1658 & $0.6846$ & 0.0443 & 0.7314 \\
Redundancy & 0.2179 & $-0.6228$ & 0.0469 & 0.6807 \\
Participation & 0.1676 & $0.3909$ & 0.0695 & 0.4609 \\
Betweenness & 0.1333 & $0.1675$ & 0.0342 & 0.2102 \\
Within-module $z$ & 0.0186 & $0.0942$ & 0.0448 & 0.1499 \\
Eigenvector & 0.0834 & $0.0926$ & 0.0317 & 0.1304 \\
k-core & 0.0556 & $0.0627$ & 0.0323 & 0.1003 \\
PageRank & 0.0379 & $0.0618$ & 0.0304 & 0.0961 \\
\bottomrule
\end{tabular}
\end{table}

Every observed association lies \textbf{inside} its null p95 band: the
pre-registered \textbf{Outcome C}. H1 is therefore \textbf{not licensed} by
this design.

\subsection{Biology (not spatially controlled) and operational chokepoints}
Top-10\%-CIS nodes under-represent cortex (0.444 vs 0.877; $q=0.00017$)
and over-represent subcortex and thalamus; these enrichments are
permutation-supported but \emph{not} spatially corrected (parcel centroids
unavailable; none invented). Cell-type fields: not available. Chokepoint
candidates (CIS top-$k$ $\wedge$ residual top-$k$ $\wedge$ bridge upper
half): 1/2/3/7 at $k=1/2/5/10\%$; node-level null verdict not established.

\subsection{Cross-scale comparison}
ME.131 (case study only): CIS $z=+23.4$ (154$\times$ peers). The fly
population-level residual fails its degree-preserving null ($z=1.21$,
$p=0.109$; frozen v1.0.0, preserved verbatim). Two same-direction
negatives license no universal-law claim.

\section{Discussion}
The human connectome's degree-independent control-impact residual is
reproducible, positive in all 801 subjects, anatomically biased away from
cortex, and substantially predictable from measured network features ---
yet its feature associations are fully matched by exact degree-preserving
nulls. The most defensible reading: what looks like a ``chokepoint
signature'' is largely a consequence of degree architecture plus
randomized-topology baselines, and any true fine structure is finer than
this design can arbitrate at parcel resolution. This bounds Paper~1
honestly: real, but not yet mechanistically explained.

\section{Limitations}
(1)~Spatial controls not established (no parcel centroids; the enrichment
is not spatially corrected). (2)~Parcel resolution: no cell-type ground
truth. (3)~Observational design: no causal claims. (4)~12-subject null
ensemble is compute-bound (bands are Monte-Carlo estimates at frozen seed;
see calibration paragraph). (5)~MLP arm unstable; reported, not
interpreted. (6)~Fly case study ($n=1$) does not generalize.

\section{Data and code availability}\label{sec:avail}
Human: AOMIC ID1000 derived connectomes (Zenodo 19796783, CC-BY-4.0)
\cite{snoek2021}; atlas: 4S456 frozen labels. Fly: FAFB v783
\cite{dorkenwald2024}. Code and artifacts:
\texttt{github.com/harsha-vardhan-2006/humanbrain-hidden-chokepoint-mechanics},
tag \texttt{v2.0.1}; the 1{,}200-record raw null ensemble ships in the
release archive (SHA256 in \texttt{dist/SHA256SUMS\_paper2.txt}).
Every stage is scripted with frozen seeds; V01--V12 re-derive all headline
numbers independently (12/12 PASS).

\section{References}
\begin{thebibliography}{99}\small
\bibitem{achard2007} Achard, S., Bullmore, E. (2007). Efficiency and cost of economical brain functional networks. \emph{PLoS Comput. Biol.} 3, e17.
\bibitem{albert2000} Albert, R., Jeong, H., Barab\'asi, A.-L. (2000). Error and attack tolerance of complex networks. \emph{Nature} 406, 378--382.
\bibitem{benjamini1995} Benjamini, Y., Hochberg, Y. (1995). Controlling the false discovery rate. \emph{J.R. Stat. Soc. B} 57, 289--300.
\bibitem{bonacich1972} Bonacich, P. (1972). Factoring and weighting approaches to status scores and clique identification. \emph{J. Math. Sociol.} 2, 113--120.
\bibitem{cliff1993} Cliff, N. (1993). Dominance statistics: ordinal analyses to answer ordinal questions. \emph{Psychol. Bull.} 114, 494--509.
\bibitem{dorkenwald2023} Dorkenwald, S. et al. (2023). Neuronal wiring diagram of an adult brain. \emph{bioRxiv} 2023.06.27.546656.
\bibitem{dorkenwald2024} Dorkenwald, S. et al. (2024). Neuronal wiring diagram of an adult brain. \emph{Nature} 634, 124--138.
\bibitem{freeman1977} Freeman, L.C. (1977). A set of measures of centrality based on betweenness. \emph{Sociometry} 40, 35--41.
\bibitem{gu2015} Gu, S. et al. (2015). Controllability of structural brain networks. \emph{Nat. Commun.} 6, 8414.
\bibitem{guimera2005} Guimer\`a, R., Amaral, L.A.N. (2005). Cartography of complex networks: modules and universal roles. \emph{J. Stat. Mech.} P02001.
\bibitem{hagmann2008} Hagmann, P. et al. (2008). Mapping the structural core of human cerebral cortex. \emph{PLoS Biol.} 6, e159.
\bibitem{hastie2009} Hastie, T., Tibshirani, R., Friedman, J. (2009). \emph{The Elements of Statistical Learning}, 2nd ed. Springer.
\bibitem{latora2001} Latora, V., Marchiori, M. (2001). Efficient behavior of small-world networks. \emph{Phys. Rev. Lett.} 87, 198701.
\bibitem{liu2011} Liu, Y.-Y., Slotine, J.-J., Barab\'asi, A.-L. (2011). Controllability of complex networks. \emph{Nature} 473, 167--173.
\bibitem{mackinnon1985} MacKinnon, J.G., White, H. (1985). Some heteroskedasticity-consistent covariance matrix estimators with improved finite sample properties. \emph{J. Econometrics} 29, 305--325.
\bibitem{maslov2002} Maslov, S., Sneppen, K. (2002). Specificity and stability in topology of protein networks. \emph{Science} 296, 910--913.
\bibitem{page1999} Page, L., Brin, S., Motwani, R., Winograd, T. (1999). The PageRank citation ranking. Stanford InfoLab technical report.
\bibitem{paper1} Malipeddi, H.V. (2026). Control-impact architecture of the human structural connectome. Paper 1 of the series, frozen release v1.1.0 (preprint / research repository).
\bibitem{seidman1983} Seidman, S.B. (1983). Network structure and minimum degree. \emph{Soc. Networks} 5, 269--287.
\bibitem{smith2015} Smith, R.E. et al. (2015). SIFT2: Super-resolution tractography. \emph{NeuroImage} 113, 323--337.
\bibitem{snoek2021} Snoek, L. et al. (2021). The Amsterdam Open MRI Collection. \emph{Sci. Data} 8, 85.
\bibitem{tang2018} Tang, E., Bassett, D.S. (2018). Colloquium: Control of networks in the brain. \emph{Rev. Mod. Phys.} 90, 031003.
\bibitem{watts1998} Watts, D.J., Strogatz, S.H. (1998). Collective dynamics of `small-world' networks. \emph{Nature} 393, 440--442.
\end{thebibliography}

\end{document}
